%%%%%%%%%%%%%%%%%%%%%%%%%%%%%%%%%%%%%%%%%
% Wenneker Assignment
% LaTeX Template
% Version 2.0 (12/1/2019)
%
% This template originates from:
% http://www.LaTeXTemplates.com
%
% Authors:
% Vel (vel@LaTeXTemplates.com)
% Frits Wenneker
%
% License:
% CC BY-NC-SA 3.0 (http://creativecommons.org/licenses/by-nc-sa/3.0/)
% 
%%%%%%%%%%%%%%%%%%%%%%%%%%%%%%%%%%%%%%%%%

%----------------------------------------------------------------------------------------
%	PACKAGES AND OTHER DOCUMENT CONFIGURATIONS
%----------------------------------------------------------------------------------------

\documentclass[11pt]{scrartcl} % Font size

%%%%%%%%%%%%%%%%%%%%%%%%%%%%%%%%%%%%%%%%%
% Wenneker Assignment
% Structure Specification File
% Version 2.0 (12/1/2019)
%
% This template originates from:
% http://www.LaTeXTemplates.com
%
% Authors:
% Vel (vel@LaTeXTemplates.com)
% Frits Wenneker
%
% License:
% CC BY-NC-SA 3.0 (http://creativecommons.org/licenses/by-nc-sa/3.0/)
% 
%%%%%%%%%%%%%%%%%%%%%%%%%%%%%%%%%%%%%%%%%

%----------------------------------------------------------------------------------------
%	PACKAGES AND OTHER DOCUMENT CONFIGURATIONS
%----------------------------------------------------------------------------------------

\usepackage{amsmath, amsfonts, amsthm} % Math packages
\usepackage{listings} % Code listings, with syntax highlighting
\usepackage[english]{babel} % English language hyphenation

\usepackage{graphicx} % Required for inserting images
\graphicspath{{Figures/}{./}} % Specifies where to look for included images (trailing slash required)

\usepackage{booktabs} % Required for better horizontal rules in tables

% \numberwithin{equation}{section} % Number equations within sections (i.e. 1.1, 1.2, 2.1, 2.2 instead of 1, 2, 3, 4)
\numberwithin{figure}{section} % Number figures within sections (i.e. 1.1, 1.2, 2.1, 2.2 instead of 1, 2, 3, 4)
\numberwithin{table}{section} % Number tables within sections (i.e. 1.1, 1.2, 2.1, 2.2 instead of 1, 2, 3, 4)

\setlength\parindent{0pt} % Removes all indentation from paragraphs

\usepackage{enumitem} % Required for list customisation
\setlist{noitemsep} % No spacing between list items

%----------------------------------------------------------------------------------------
%	DOCUMENT MARGINS
%----------------------------------------------------------------------------------------

\usepackage{geometry} % Required for adjusting page dimensions and margins

\geometry{
	paper=a4paper, % Paper size, change to letterpaper for US letter size
	top=2.5cm, % Top margin
	bottom=3cm, % Bottom margin
	left=3cm, % Left margin
	right=3cm, % Right margin
	headheight=0.75cm, % Header height
	footskip=1.5cm, % Space from the bottom margin to the baseline of the footer
	headsep=0.75cm, % Space from the top margin to the baseline of the header
	%showframe, % Uncomment to show how the type block is set on the page
}

%----------------------------------------------------------------------------------------
%	FONTS
%----------------------------------------------------------------------------------------

\usepackage[utf8]{inputenc} % Required for inputting international characters
\usepackage[T1]{fontenc} % Use 8-bit encoding

\usepackage{fourier} % Use the Adobe Utopia font for the document

%----------------------------------------------------------------------------------------
%	SECTION TITLES
%----------------------------------------------------------------------------------------

\usepackage{sectsty} % Allows customising section commands
\sectionfont{\vspace{6pt}\normalfont\scshape} % \section{} styling
\subsectionfont{\normalfont\bfseries} % \subsection{} styling
\subsubsectionfont{\normalfont\itshape} % \subsubsection{} styling
\paragraphfont{\normalfont\scshape} % \paragraph{} styling

%----------------------------------------------------------------------------------------
%	HEADERS AND FOOTERS
%----------------------------------------------------------------------------------------

\usepackage{scrlayer-scrpage} % Required for customising headers and footers

\ohead*{} % Right header
\ihead*{} % Left header
\chead*{} % Centre header

\ofoot*{} % Right footer
\ifoot*{} % Left footer
\cfoot*{\pagemark} % Centre footer
 % Include the file specifying the document structure and custom commands

%----------------------------------------------------------------------------------------
%	TITLE SECTION
%----------------------------------------------------------------------------------------
% Add a logo to the title page from the folder "Figures"





\title{	
	\includegraphics[width=0.25\textwidth]{./Figures/Sharif_University_Logo.jpg}\\
	\normalfont\normalsize
	\textsc{Sharif University of Technology}\\ % Your university, school and/or department name(s)
	\vspace{25pt} % Whitespace
	\rule{\linewidth}{0.5pt}\\ % Thin top horizontal rule
	\vspace{20pt} % Whitespace
	{\huge Deep Learning Assignment 2}\\ % The assignment title
	\vspace{12pt} % Whitespace
	\rule{\linewidth}{2pt}\\ % Thick bottom horizontal rule
	\vspace{12pt} % Whitespace
}
% \includegraphics[width=0.1\textwidth]{C:/Users/saeedzou/Documents/sharif_logo.png}~[1cm]

\author{\LARGE Saeedreza Zouashkiani} % Your name
\DeclareMathOperator{\E}{\mathbb{E}}
\DeclareMathOperator*{\argmin}{argmin}

\date{\normalsize\today} % Today's date (\today) or a custom date

\begin{document}
% \titlegraphic{\includegraphics[width=0.5\textwidth]{./Figures/sharif_logo.png}}
\maketitle % Print the title

\section{} %1


	\subsection{}
		Consider the Taylor expansion of the loss function at $\boldsymbol\omega$:
		\begin{equation} 
			\boldsymbol{J}(\boldsymbol{\omega}+ \delta\boldsymbol\omega) - \\
			\boldsymbol{J}(\boldsymbol{\omega}) \approx \\
			\nabla_{\boldsymbol\omega}\boldsymbol{L}\delta\boldsymbol{\omega} + \\
			\frac{1}{2}\delta\boldsymbol{\omega}^T\boldsymbol{H}\delta\boldsymbol{\omega} 
		\end{equation}

		The first term goes to zero when the algorithm converges. So, our problem is to minimize:
		\begin{equation}\label{eq:1}
			\boldsymbol{C} = \frac{1}{2}\delta\boldsymbol{\omega}^T\boldsymbol{H}\delta\boldsymbol{\omega}
		\end{equation}

		Our goal is then to set one of the weights to 
		zero which is denoted by $\omega_i$ to minimize the increase in error given by equation (\ref{eq:1}). 
		Then to eliminate $\omega_i$:
		\begin{equation}\label{eq:2}
			\omega_i + \boldsymbol{e_i}^T.\delta\boldsymbol\omega = 0
		\end{equation}
		Where $\boldsymbol{e_i}$ is a vector with a 1 corresponding to $\omega_i$ and zero elsewhere.
		To minimize equation (\ref{eq:1}) constrainted on equation (\ref{eq:2}), we use Lagrangian multipliers:
		\begin{equation}
			\boldsymbol{L} = \frac{1}{2}\delta\boldsymbol{\omega}^T\boldsymbol{H}\delta\boldsymbol{\omega} \\
			+ \lambda(\omega_i + \boldsymbol{e_i}^T.\delta\boldsymbol\omega)
		\end{equation}
		Differentiating with respect to $\delta\boldsymbol\omega$ and setting to zero:
		\begin{equation} \label{eq:3}
			\boldsymbol{H}\delta\boldsymbol{\omega} + \lambda \boldsymbol{e_i}=0
		\end{equation}
		Then if we solve for $\lambda$:
		\begin{equation}\label{eq:4}
			\lambda = \frac{\omega_i}{(\boldsymbol{H^{-1}})_{ii}}
		\end{equation}
		Where $\boldsymbol{H^{-1}}$, corresponds to the $i^{th}$ element in the inverse Hessian matrix.
		Then by plugging (\ref{eq:4}) into (\ref{eq:3})
		\begin{equation} \label{eq:5}
			\delta\boldsymbol{\omega} = -\frac{\omega_i}{(\boldsymbol{H^{-1}})_{ii}}\\
			 \boldsymbol{H^{-1}}.\boldsymbol{e_i}
		\end{equation}
		Then plugging (\ref{eq:5}) into (\ref{eq:1}):
		\begin{equation}
			\boldsymbol{C} = \frac{1}{2} \frac{\omega_i^2}{(\boldsymbol{H^{-1}})_{ii}}
		\end{equation}

	\subsection{}
		If $ \boldsymbol{H} = \boldsymbol{I} $, then $ \boldsymbol{H^{-1}} = \boldsymbol{I} $. Therefore:
		\begin{equation}
			\delta\boldsymbol{\omega} = -\omega_i
		\end{equation}
		\begin{equation}
			\boldsymbol{C} = \frac{1}{2} \omega_i^2
		\end{equation}
		In the previous part, we needed to compute the inverse Hessian first, and compute 
		$\frac{1}{2} \frac{\omega_i^2}{(\boldsymbol{H^{-1}})_{ii}}$ for each element and remove 
		the one with the
		least increase in the loss, but in here we compute $\frac{1}{2} \omega_i^2$ 
		and set $\delta\boldsymbol{\omega} = -\omega_i$
		for that $i$ that minimizes the loss.
\section{} %2
		\subsection{}
		This is solved in part 3.1
		\subsection{}
		\begin{equation}
			\begin{aligned} \label{eq:2.2}
				&\mathbb{E}[\frac{1}{N} ||\boldsymbol{ (X(X^T X)^{-1}X^T - I)\epsilon }||_2^2]  \\
				&=\frac{1}{N} \mathbb{E}[\boldsymbol{\epsilon^T (X(X^T X)^{-1}X^T - I)^T (X(X^T X)^{-1}X^T - I) \epsilon }]  \\
				&=\frac{1}{N} \mathbb{E}[\boldsymbol{\epsilon^T (B)^T (B) \epsilon }]
			\end{aligned}
		\end{equation}
		Since $\boldsymbol{A}$ is symmetric and $\boldsymbol{A^T A}=\boldsymbol{A^2}$, then it is idempotent. That is
		its eigenvalues consist of $d$ ones and $n-d$ zeros. Therefore using \ref{eq:2.2}:
		\begin{equation}
			\boldsymbol{B^T B = (A - I)^T (A - I) = A^2 - 2A + I = I - A }  
		\end{equation}
		Since $\boldsymbol{\epsilon^T B^T B \epsilon }$ is a scalar, then its trace equals itself:
		\begin{equation}
			\begin{aligned}
				&\frac{1}{N} \mathbb{E}[\boldsymbol{\epsilon^T B^T B \epsilon }]  \\
				&=\frac{1}{N} \mathbb{E}[\text{Trace}(\boldsymbol{\epsilon^T B^T B \epsilon })]  \\
				&=\frac{1}{N} \mathbb{E}[\text{Trace}(\boldsymbol{ B^T B \epsilon^T \epsilon })]  \\
				&=\frac{1}{N} \text{Trace}(\boldsymbol{ B^T B \Sigma})  \\
				&=\frac{\sigma^2}{N} \text{Trace}(\boldsymbol{I - A})  \\
				&=\frac{N-d}{N} \sigma^2
			\end{aligned}
		\end{equation}
		Where the last part comes from the propertiy of idempotent matrices. That is $\text{Trace}(\boldsymbol{A}) = \text{rank}(\boldsymbol{X})$.

		\subsection{}
		As d increases and gets closer to N, the trace of $\boldsymbol{I - A}$ approaches zero. That is 
		the number of independent columns of $\boldsymbol{X}$ increases, so we can have a better projection matrix than before.
		Naturally the training error will decrease as well.

			
\section{} %3
	\subsection{} %3.1
		We want to find:
		\begin{equation}
			\begin{aligned}
				&\argmin_{\boldsymbol{\beta}} ||\boldsymbol{Y}-\boldsymbol{X\beta}||_2^2 \\
				&=\argmin_{\boldsymbol{\beta}} (\boldsymbol{Y}-\boldsymbol{X\beta})^T (\boldsymbol{Y}-\boldsymbol{X\beta})\\
				&=\argmin_{\boldsymbol{\beta}} \boldsymbol{Y}^T\boldsymbol{Y} - 2\boldsymbol{\beta}^T\boldsymbol{X}^T\boldsymbol{Y} +
				\boldsymbol{\beta}^T\boldsymbol{X}^T\boldsymbol{X\beta}\\
				&=\argmin_{\boldsymbol{\beta}}  - 2\boldsymbol{\beta}^T\boldsymbol{X}^T\boldsymbol{Y} +
				\boldsymbol{\beta}^T\boldsymbol{X}^T\boldsymbol{X\beta}
			\end{aligned}
		\end{equation}
		Define $L(\boldsymbol{\beta})=2\boldsymbol{\beta}^T\boldsymbol{X}^T\boldsymbol{Y} +
		\boldsymbol{\beta}^T\boldsymbol{X}^T\boldsymbol{X\beta}$, and differentiate with respect to $\boldsymbol{\beta}$
		and set to zero:
		\begin{equation}
			\frac{\partial{L(\boldsymbol{\beta})}}{\boldsymbol{\beta}}=
			-2\boldsymbol{X}^T\boldsymbol{Y} + 2\boldsymbol{X}^T\boldsymbol{X}\boldsymbol{\beta}
		\end{equation}
		\begin{equation}
			\boldsymbol{\hat{\beta}} = (\boldsymbol{X}^T\boldsymbol{X})^{-1}\boldsymbol{X}^T\boldsymbol{Y}
		\end{equation}
	\subsection{} %3.2
		Add the L2 norm of the weights to the minimization problem of the previous part:
		\begin{equation}
			\begin{aligned}
				&\argmin_{\boldsymbol{\beta}} ||\boldsymbol{Y}-\boldsymbol{X\beta}||_2^2 
				+\alpha||\boldsymbol{\beta}||_2^2\\
				&=\argmin_{\boldsymbol{\beta}} (\boldsymbol{Y}-\boldsymbol{X\beta})^T (\boldsymbol{Y}-\boldsymbol{X\beta})+
				\alpha\boldsymbol{\beta}^T\boldsymbol{\beta}\\
				&=\argmin_{\boldsymbol{\beta}} \boldsymbol{Y}^T\boldsymbol{Y} - 2\boldsymbol{\beta}^T\boldsymbol{X}^T\boldsymbol{Y} +
				\boldsymbol{\beta}^T\boldsymbol{X}^T\boldsymbol{X\beta}+
				\alpha\boldsymbol{\beta}^T\boldsymbol{\beta}\\
				&=\argmin_{\boldsymbol{\beta}}  - 2\boldsymbol{\beta}^T\boldsymbol{X}^T\boldsymbol{Y} +
				\boldsymbol{\beta}^T\boldsymbol{X}^T\boldsymbol{X\beta}+
				\alpha\boldsymbol{\beta}^T\boldsymbol{\beta}
			\end{aligned}
		\end{equation}
		Define $L(\boldsymbol{\beta})=2\boldsymbol{\beta}^T\boldsymbol{X}^T\boldsymbol{Y} +
		\boldsymbol{\beta}^T\boldsymbol{X}^T\boldsymbol{X\beta}
		+\alpha\boldsymbol{\beta}^T\boldsymbol{\beta}$
		, and differentiate with respect to $\boldsymbol{\beta}$
		and set to zero:
		\begin{equation}
			\frac{\partial{L(\boldsymbol{\beta})}}{\boldsymbol{\beta}}=
			-2\boldsymbol{X}^T\boldsymbol{Y} + 2\boldsymbol{X}^T\boldsymbol{X}\boldsymbol{\beta}+
			2\alpha\boldsymbol{\beta}
		\end{equation}
		\begin{equation}
			\boldsymbol{\hat{\beta}} = (\boldsymbol{X}^T\boldsymbol{X} 
			+ \alpha\boldsymbol{I})^{-1}\boldsymbol{X}^T\boldsymbol{Y}
		\end{equation}

	\subsection{} %3.3
		We know that both $\boldsymbol{\beta}^*$ and $\boldsymbol{\hat{\beta}}$ have an expectation value of $\boldsymbol{\beta}$.
		Let $\boldsymbol{\beta}^* = \boldsymbol{Ay}$ and $\boldsymbol{\hat{\beta}}= \boldsymbol{By}$.
		Where $\boldsymbol{A} = (\boldsymbol{X}^T\boldsymbol{\Sigma^{-1}X})^{-1}\boldsymbol{\Sigma^{-1}X^T}$ 
		and $\boldsymbol{B} = (\boldsymbol{X}^T\boldsymbol{X})^{-1}\boldsymbol{X}^T$.
		\begin{equation}
			\begin{aligned}
				Var(\boldsymbol{\beta}^*) &=  \boldsymbol{B}Var(\boldsymbol{y})\boldsymbol{B}^T = \boldsymbol{X^T \Sigma^{-1} X}^{-1}\\
				Var(\boldsymbol{\hat{\beta}}) &= \boldsymbol{A}Var(\boldsymbol{y})\boldsymbol{A}^T =
				\boldsymbol{(X^T X)^{-1} X^T \Sigma X (X^T X)^{-1}}
			\end{aligned}
		\end{equation}
		If $\boldsymbol{XF}=\boldsymbol{\Sigma X}$, then:
		\begin{equation} \label{eq:3.3}
		 (\boldsymbol{X^T \Sigma^{-1} X})^{-1} = \boldsymbol{F(X^T X)^{-1}}
		\end{equation}
		It is obvious from \ref{eq:3.3} that $\boldsymbol{F}$ is non-singular. Using (\ref{eq:3.3}), we can write:
		\begin{equation}
			\begin{aligned}
				Var(\boldsymbol{\beta}^*) & = \boldsymbol{F(X^T X)^{-1}} \\
				Var(\boldsymbol{\hat{\beta}}) & =
				\boldsymbol{(X^T X)^{-1} X^T \Sigma X (X^T X)^{-1}} =
				\boldsymbol{(X^T X)^{-1} X^T X F (X^T X)^{-1}} = 
				\boldsymbol{F(X^T X)^{-1}}
			\end{aligned}
		\end{equation}
		Therefore $\boldsymbol{\beta}^*$ and $\boldsymbol{\hat{\beta}}$ estimate the same quantity.
		On the other hand, to prove when the estimators are equal, we need to show that:
		\begin{equation}
			\begin{aligned}
				\boldsymbol{\beta}^* &= \boldsymbol{\hat{\beta}} \\
				(\boldsymbol{X^T \Sigma^{-1} X})^{-1} &= \boldsymbol{(X^T X)^{-1} X^T \Sigma X (X^T X)^{-1}} \\
				\boldsymbol{X^T X}(\boldsymbol{X^T \Sigma^{-1} X})^{-1} &= \boldsymbol{X^T \Sigma X (X^T X)^{-1}} \\
				\boldsymbol{X^T} \left(\boldsymbol{X}\boldsymbol{(X^T \Sigma^{-1} X)^{-1} - \Sigma X (X^T X)^{-1}}\right)&=0\\
				\boldsymbol{X}\boldsymbol{(X^T \Sigma^{-1} X)^{-1}} &= \boldsymbol{\Sigma X (X^T X)^{-1}}\\
				\boldsymbol{X}\boldsymbol{(X^T \Sigma^{-1} X)^{-1} X^T X} &= \boldsymbol{\Sigma X}\\
				\boldsymbol{XF} &= \boldsymbol{\Sigma X} 
			\end{aligned}
		\end{equation}
	\subsection{} 
	Extend the error function without the L-1 norm:
	\begin{equation} \label{eq:3.4}
		\begin{aligned}
			L(\boldsymbol{\beta}, \lambda_1) &= 
			||\boldsymbol{Y}-\boldsymbol{X\beta}||_2^2 + \lambda_1||\boldsymbol{\beta}||_2^2=
			(\boldsymbol{Y}-\boldsymbol{X\beta})^T (\boldsymbol{Y}-\boldsymbol{X\beta}) +
			\lambda_1 \boldsymbol{\beta}^T \boldsymbol{\beta} \\
			&=\boldsymbol{Y}^T\boldsymbol{Y} - 2\boldsymbol{\beta}^T\boldsymbol{X}^T\boldsymbol{Y} +
			\boldsymbol{\beta}^T\boldsymbol{X}^T\boldsymbol{X\beta}+
			\lambda_1 \boldsymbol{\beta}^T \boldsymbol{\beta}
		\end{aligned}
	\end{equation}
	Define the augmented $\boldsymbol{\hat{X}} = \begin{bmatrix}
		\boldsymbol{X}\\
		\sqrt{\lambda_1}\boldsymbol{I}
	\end{bmatrix}$
	, and $\boldsymbol{\hat{Y}} = \begin{bmatrix}
		\boldsymbol{Y}\\
		\boldsymbol{0}
	\end{bmatrix}$. Then equation (\ref{eq:3.4}) can be written as:
	\begin{equation}
		L(\boldsymbol{\beta}, \lambda_1) = 
			||\boldsymbol{\hat{Y}}-\boldsymbol{\hat{X}\beta}||_2^2
	\end{equation}
	Therefore the loss function $L(\boldsymbol{\beta}, \lambda_1, \lambda_2) $ will become:
	\begin{equation}
		L(\boldsymbol{\beta}, \lambda_1, \lambda_2) = ||\boldsymbol{\hat{Y}}-\boldsymbol{\hat{X}\beta}||_2^2
	\end{equation}


\section{} %4
	\subsection{}
		Let $J_D = (y_d- O_D)^2$ and $J_N = (y_d- O_N)^2$
		denote the loss function of the Dropout and without Dropout network, respectively, and 
		$O_D = \sum_{k=1}^n{\delta_k\omega_k x_k}$
		\begin{equation}
			\begin{aligned}
				\frac{\partial{J_N}}{\partial{\omega_i}} &= 
				-2(y_d - O_N)\frac{\partial{O_N}}{\partial{\omega_i}} =
				-2(y_d - O_N)x_i = 
				-2y_d x_i + 2 \sum_{k=1}^n{\omega_k x_k x_i} \\
				&= -2y_d x_i +2\omega_i x_i^2 + 2 \sum_{k=1, k\neq{i}}^n{\omega_k x_k x_i}
			\end{aligned}
		\end{equation}
		
		\begin{equation}
			\begin{aligned}
			\frac{\partial{J_D}}{\partial{\omega_i}} &= 
			-2(y_d - O_D)\frac{\partial{O_D}}{\partial{\omega_i}} =
			-2(y_d - O_D)\delta_i x_i \\
			&= -2y_d \delta_i x_i + 2\omega_i^2 \delta_i^2 x_i^2 + 2\sum_{k=1, k\neq{i}}^n{\omega_k \delta_k \delta_i x_k x_i}
		\end{aligned}
		\end{equation}
		Since equation (12) is probabilistic, we take the expectation:
		\begin{equation}
			\begin{aligned}
				\mathbb{E}[\frac{\partial{J_D}}{\partial{\omega_i}}] &=
				-2y_d x_i +2\omega_i x_i^2 + 2\sigma^2 \omega_i^2 x_i^2 + 2\sum_{k=1, k\neq{i}}^n{\omega_k x_k x_i} \\
				&= \frac{\partial{J_N}}{\partial{\omega_i}} + 2\sigma^2 \omega_i^2 x_i^2
			\end{aligned}
		\end{equation}
		\subsection{}
			The overall regularized loss function can be written as:
			\begin{equation}
				J_R = (y_d- \sum_{k=1}^n{\delta_k\omega_k x_k})^2 + \sigma^2 \sum_{k=1}^n {w_k^2 x_k^2}
			\end{equation}
			It means that the network is regularized by the square of multiplication of weights by inputs
			with a regularization factor of $\sigma^2$.
\section{} %5
	Let the cost function be:
	\begin{equation}
		L = \frac{1}{2}\boldsymbol{\omega^T}\boldsymbol{H}\boldsymbol{\omega}
	\end{equation}
	Computing the first and second order derivatives:
	\begin{equation}
		\boldsymbol{\nabla_{\omega}}L = \frac{\partial{L}}{\partial{\boldsymbol{\omega}}} =\boldsymbol{H}\boldsymbol{\omega}
	\end{equation}
	\begin{equation} \label{eq:5.1}
		\boldsymbol{\nabla^2_{\omega}}L =\frac{\partial^2{L}}{\partial{\boldsymbol{\omega^2}}} =\boldsymbol{H}
	\end{equation}
	\subsection{}
		\begin{equation}
			\boldsymbol{\omega}^{(t+1)} = \boldsymbol{\omega}^{(t)} - 
			\epsilon\boldsymbol{H}\boldsymbol{\omega}^{(t)} =
			\boldsymbol{Q}(\boldsymbol{I} - \epsilon \boldsymbol{\Lambda})\boldsymbol{Q^T \omega^{(t)}}
		\end{equation}
	\subsection{}
		\begin{equation}
			\boldsymbol{\omega}^{(t+1)} = 
			\boldsymbol{Q}(\boldsymbol{I} - \epsilon \boldsymbol{\Lambda})\boldsymbol{Q^T \omega^{(0)}}
		\end{equation}
	\subsection{}
		$|(\boldsymbol{I} - \epsilon \boldsymbol{\Lambda})|<\boldsymbol{1}$ must be satisfied. Therefore:
		\begin{equation}
			\epsilon < \frac{2}{\lambda_{max}}
		\end{equation}
	\subsection{}
		Because of \ref{eq:5.1}, the newtons method will be:
		\begin{equation}
			\boldsymbol{\omega}^{(t+1)} = (1-\epsilon)\boldsymbol{\omega}^{(t)}
		\end{equation}
		Then by choosing $\epsilon = 1$, the newton method will converge in one step.
	\subsection{}
	 Because in Newton's method, calculating the inverse of the Hessian matrix is expensive,
	 gradient descent is used instead.

\section{} %6
\subsection{}
It will learn a similarity measure between 2 inputs, since the inputs are both mapped to a latent
space and the distance between the 2 latent vectors is minimized. This can be used for images 
and text, for example.
\subsection{}
\begin{equation}
	\begin{aligned}
		\boldsymbol{z_1} &= \boldsymbol{Wx_1 + b} \\
		\boldsymbol{z_2} &= \boldsymbol{Wx_2 + b} \\
		\boldsymbol{\Delta_1 } &= \frac{\partial{\boldsymbol{J}}}{\partial{\boldsymbol{h_1}}} =
		2\boldsymbol{h_1 - h_2} \\
		\boldsymbol{\Delta_2 } &= \frac{\partial{\boldsymbol{J}}}{\partial{\boldsymbol{h_2}}} =
		2\boldsymbol{h_2 - h_1} \\
		\boldsymbol{\Delta_3} &= \frac{\partial{\boldsymbol{J}}}{\partial{\boldsymbol{z_1}}} =
		\boldsymbol{\Delta_1} \bigodot \boldsymbol{(1 - h_1^2)} \\
		\boldsymbol{\Delta_4} &= \frac{\partial{\boldsymbol{J}}}{\partial{\boldsymbol{z_2}}} =
		\boldsymbol{\Delta_2} \bigodot \boldsymbol{(1 - h_2^2)} \\
		\frac{\partial{\boldsymbol{J}}}{\partial{\boldsymbol{W}}} &= 
		\boldsymbol{\Delta_3 x_1^T +  \Delta_4 x_2^T} + 2 \boldsymbol{W} \\
		\frac{\partial{\boldsymbol{J}}}{\partial{\boldsymbol{b}}} &= 
		\boldsymbol{\Delta_3 +  \Delta_4} \\
	\end{aligned}
\end{equation}
Then the gradient descent update rule for $m$ batch size is:
\begin{equation}
	\begin{aligned}
		\boldsymbol{W}^{(t+1)} = \boldsymbol{W}^{(t)} - \epsilon \frac{1}{m}\sum_{i=1}^m\frac{\partial{\boldsymbol{J}_i}}{\partial{\boldsymbol{W}}} \\
		\boldsymbol{b}^{(t+1)} = \boldsymbol{b}^{(t)} - \epsilon \frac{1}{m}\sum_{i=1}^m\frac{\partial{\boldsymbol{J}_i}}{\partial{\boldsymbol{b}}} \\
	\end{aligned}
\end{equation}


%----------------------------------------------------------------------------------------
%	FIGURE EXAMPLE
%----------------------------------------------------------------------------------------
% \begin{figure}[h] % [h] forces the figure to be output where it is defined in the code (it suppresses floating)
% 	\centering
% 	\includegraphics[width=0.5\columnwidth]{swallow.jpg} % Example image
% 	\caption{European swallow.}
% \end{figure}

%------------------------------------------------
%----------------------------------------------------------------------------------------
%	TEXT EXAMPLE
%----------------------------------------------------------------------------------------


\end{document}
